% !TeX spellcheck = ru_RU
% !TEX root = vkr.tex

\section*{Заключение}

В данной работе был создан ML-сервис для сбора статистики со строительных объектов,
способный собирать статистику на основе анализа видео, в котором видна зона ворот.
Также в этом сервисе присутствует \texttt{Telegram}-бот, который отправляет отчёты подписанным пользователям.
В процессе работы были выполнены поставленные задачи.

\begin{enumerate}
    \item Был проведён обзор существующих решений (подробнее в разделе~\ref{sec:related}).
    \item Разработан способ сбора статистики и последующей обработки видео (подробнее в подразделе~\ref{subsec:getting_data_from_objects}).
    \item Был реализован ML-алгоритм, обрабатывающий данные с помощью дообученной модели \texttt{YOLO} (подробнее в подразделах~\ref{subsec:creating_ml_pipeline} и~\ref{subsubsec:yolo-training-results}).
    \item Разработан \texttt{Telegram}-бот для ежедневной рассылки результатов обработки (подробнее в подразделе~\ref{subsec:creating_system_arch}).
    \item Разработан формат статистического отчёта для рассылки (подробнее в подразделе~\ref{subsec:creating_report}).
    \item Получены метрики обучения модели и работы end-to-end сервиса. Реализация может использоваться для дальнейших проектов (подробнее в подразделе~\ref{subsec:experiment-results}).
\end{enumerate}

Исходный код данной работы: \url{https://github.com/VadosikRRR/ObjectStatistick}
